\section{Introduction}
A \emph{real zero polynomial} is a polynomial $P\in\R[x_1,\ldots,x_n]$ with $P(0)\neq0$ which has only real zeros when restricted to any real line through the origin. Real zero polynomials play an important role in the theory of semidefinite programming, see \eg \cite{victorsurvey}. In this context, the second author and Schweighofer made the following conjecture:

\begin{con*}[{\cite[Conjecture~7.6]{rzamalgamation}}]
    Assume that $F\in\R[x_1,x_2,y_1,\ldots,y_m]$ and $G\in\R[x_1,x_2,z_1,\ldots,z_n]$ are real zero polynomials. If
    \begin{equation*}
        F|_{y_1=\cdots=y_m=0}=G|_{z_1=\cdots=z_n=0},
    \end{equation*}
    then there is a real zero polynomial $H\in\R[x_1,x_2,y_1,\ldots,y_m,z_1,\ldots,z_n]$ such that
    \begin{equation*}
        F=H|_{z_1=\cdots=z_n=0}\textrm{ and }G=H|_{y_1=\cdots=y_m=0}.
    \end{equation*}
\end{con*}

We provide a counterexample to this conjecture in \Cref{sec:amal}. Using the connection between stable polynomials, a concept closely related to real zero polynomials, and the theory of matroids and polymatroids established in \cite{HPPlong} and \cite{BrandenHPP}, we associate a polymatroid to every real zero polynomial. Then the key step is \Cref{thm:amalmain} which says that if two real zero polynomials satisfy the conclusion the above conjecture, then the associated polymatroids can be \emph{amalgamated} in the sense of \Cref{def:amal}. Our counterexample comes from two polymatroids which cannot be amalgamated. These polymatroids are obtained by specializing the matroids $F_7^{-4}$ and $F_7^{-5}$ --- two relaxations of the Fano matroid which have the \emph{half-plane property}, meaning that their bases generating polynomials are stable. 

More generally, a matroid has the \emph{weak half-plane property} if its set of bases agrees with the support of a stable polynomial. The question of which matroids have the (weak) half-plane property has been extensively studied. In this context, Br\"and\'en and D'Le\'on offered the following conjecture:

\begin{con*}[{\cite[Conjecture~4.2]{BrandenDLeon}}]
    Suppose that $M$ has the weak half-plane property. Then so does any relaxation of $M$.
\end{con*}

In \Cref{sec:whpp} we present a counterexample to this conjecture. More precisely, we show that a relaxation $P_1$ of the matroid $P_8$ does not have the weak half-plane property. On the other hand, $P_8$ does have the weak half-plane property because it is representable over $\R$. For proving that $P_1$ does not have the weak half-plane property, we employ the techniques developed in \cite{BrandenDLeon} for narrowing down the space of possible coefficients of a hypothetical stable polynomial with support $P_1$.

A key feature of the matroid $P_1$ that made us examine it for the weak half-plane property is that it is not representable --- even under more general notions of representability as discussed in \cite{skewpartial}. Matroids representable over more general structures than fields still often tend to have the weak half-plane property, see \eg \cite{NonPappus}. Our third result is of this flavor. Namely, we prove in \Cref{sec:qu} a conjecture which was attributed by Pendavingh and van Zwam to David G. Wagner:

\begin{con*}[{\cite[Conjecture~6.9]{skewpartial}}]
    All quaternionic unimodular matroids have the half-plane property.
\end{con*}

The notion of quaternionic unimodular matroids is a generalization of the class of sixth root of unity matroids to the skew field of quaternions. For a precise definition see \Cref{def:qu}.

\section{Preliminaries}
We denote by $\N$ and $\N_0$ the set of positive and nonnegative integers, respectively.
In this section we let $E$ always denote a finite set. For $i\in E$ we denote by $\delta_i\in\R^E$ the $i$th unit vector. For $x\in\R^E$ we write $|x|=\sum_{i\in E}|x_i|$. We further let $[n]=\{1,\ldots,n\}$ and denote by $\binom{[n]}{k}$ the set of all $k$-element subsets of $[n]$ for every $k,n\in\N$. We recall the cryptomorphic definitions of M-convex sets and polymatroids.

\begin{Def}
     A subset $J\subseteq\N_0^E$ is \emph{M-convex} if for every $i\in E$ and every $\alpha,\beta\in J$ such that $\alpha_i>\beta_i$, there is $j\in E$ satisfying
            \begin{equation*}
                \alpha_j<\beta_j\textrm{ and }\alpha-\delta_i+\delta_j\in J.
            \end{equation*}
\end{Def}

\begin{Def}
    A \emph{polymatroid} on $E$ is a function $r\colon 2^E\to\N_0$ such that we have for all $S,T\subseteq E$:
            \begin{enumerate}[(i)]
            \item $r(\emptyset)=0$,
                \item $r(S)\leq r(T)$ if $S\subseteq T$, and
                \item $r(S\cup T)+r(S\cap T)\leq r(S)+r(T)$.
            \end{enumerate}
\end{Def}
If $J\subseteq\N_0^E$ is an M-convex set, then we define the function $r_J\colon2^E\to\N_0$ by
    \begin{equation*}
        r_J(S)=\max\{\sum_{i\in S}\alpha_i\mid \alpha\in J\}.
    \end{equation*}
Conversely, if $r\colon 2^E\to\N_0$ is a polymatroid, then we define the set $J_r\subseteq\N_0^E$ by
\begin{equation*}
    J_r=\{x\in\N_0^E\mid \sum_{i\in S}x_i\leq r(S)\textrm{ for all }S\subseteq E\textrm{ and }\sum_{i\in E}x_i=r(E)\}.
\end{equation*}
These two constructions are inverse to each other and define bijections between the set of M-convex sets in $\N_0^E$ and the set of polymatroids on $E$, see for example \cite[\S4.4]{murota}.

\begin{rem}
    A polymatroid $r\colon 2^E\to\N_0$ is the rank function of a matroid $M$ on $E$ if and only if $r(\{i\})\leq1$ for all $i\in E$. In this case we have
    \begin{equation*}
        J_r=\{\sum_{i\in B}\delta_i\mid B\textrm{ is a basis of }M\}.
    \end{equation*}
    We apply the definitions made for polymatroids to matroids by considering their rank functions as polymatroids.
\end{rem}

\subsection{Amalgamation of polymatroids}
There has been considerable interest in the question whether two (poly)matroids can be amalgamated in the following sense, see for example \cite{sticky,stickypolymatroids} and also \cite[\S11.4]{oxley}.
\begin{Def}\label{def:amal}
    Let $E_1,E_2$ be two finite sets, let $E_0=E_1\cap E_2$ and $E_3=E_1\cup E_2$. Let $r_i$ be a polymatroid on $E_i$ for all $i=0,1,2,3$.
    \begin{enumerate}[(a)]
        \item We say that $r_0$ is the \emph{restriction} of $r_1$ to $E_0$ if $r_0(S)=r_1(S)$ for all $S\subseteq E_0$ and denote this by $r_0=r_1|_{E_0}$.
        \item We say that $r_3$ is an \emph{amalgam} of $r_1$ and $r_2$ if $r_1=r_3|_{E_1}$ and $r_2=r_3|_{E_2}$.
    \end{enumerate}
\end{Def}

Clearly, a necessary condition for an amalgam of $r_1$ and $r_2$ as in \Cref{def:amal} to exist is that $r_1|_{E_0}=r_2|_{E_0}$. However, this condition is not sufficient \cite[Example~7.2.4]{oxley}.

\begin{Def}
    A polymatroid $r_0$ on $E$ is called \emph{sticky} if there exists an amalgam for all polymatroids $r_1$ and $r_2$ on finite sets $E_1$ and $E_2$ with $E=E_1\cap E_2$ and such that $r_0=r_i|_{E}$ for $i=1,2.$
\end{Def}

In order to describe conditions for a polymatroid to be sticky, recall the following definition, generalizing the usual notion for matroids.
\begin{Def}
    Let $r$ be a polymatroid on $E$.
    \begin{enumerate}[(a)]
        \item A subset $F\subseteq E$ is called a \emph{flat} of $r$ if $r(F')>r(F)$ for every proper superset $F'$ of $F$.
        \item Two flats $F_1,F_2$ are called a \emph{modular pair} if
        \begin{equation*}
            r(F_1)+r(F_2)=r(F_1\cup F_2)+r(F_1\cap F_2).
        \end{equation*}
    \end{enumerate}
\end{Def}

\begin{thm}[\cite{stickypolymatroids}]\label{thm:stickypoly}
    Let $r$ be a polymatroid on $E$. If every pair of flats is modular, then $r$ is sticky. The converse holds if $|E|\leq5$.
\end{thm}

\begin{ex}\label{ex:uniform}
 Let $2\leq n\leq5$ and $E=[n]$. One checks that
 \begin{equation*}
     J=\{x\in\N_0^n\mid |x|=3\textrm{ and }\forall i\in E: x_i\leq 2\}
 \end{equation*}
 is M-convex. The polymatroid $r_J$ is not sticky because the pair $\{1\},\{2\}$ of flats is not modular.
\end{ex}

\subsection{Stable polynomials}
We briefly recall the different stability and real zero properties of polynomials and their relation to each other. As a general reference we recommend the survey \cite{pemantle}.
\begin{Def}
    Let $P\in\R[x_i\mid i\in E]$ be a polynomial.
    \begin{enumerate}[(a)]
        \item The \emph{support} of $P$ is the unique subset $\supp(P)\subseteq\N_0^E$ such that we can write
        \begin{equation*}
            P=\sum_{\alpha\in\supp(P)}c_\alpha x^\alpha
        \end{equation*}
        for some non-zero $c_\alpha\in\R$.
        \item The polynomial $P$ is called \emph{stable} if for all $z\in\C^E$ such that $\Ima(z_i)>0$ for all $i\in E$ we have $P(z)\neq0$.
        \item The polynomial $P$ is called a \emph{real zero polynomial} if for all $v\in\R^E$ the univariate polynomial $P(t\cdot v)\in\R[t]$ has only real zeros.
        \item Let $P$ be homogeneous and $e\in\R^E$. Then $P$ is called \emph{hyperbolic} with respect to $e$ if for all $v\in\R^E$ the univariate polynomial $P(t\cdot e+v)\in\R[t]$ has only real zeros.
    \end{enumerate}
\end{Def}

Taking $v=0$ in the definition of a real zero polynomial, yields that a real zero polynomial does not vanish at the origin. Similarly, a polynomial $P$ which is hyperbolic in direction $e$, satisfies $P(e)\neq0$. Stability is preserved under taking partial derivatives in coordinate directions, under setting some variables equal to each other and under scaling variables by positive scalars. We summarize the well-known connection between the above concepts.

\begin{prop}[see for example {\cite[Proposition 5.3]{pemantle}}]\label{prop:stablehyp}
    Let $P\in\R[x_i\mid i\in E]$ be a homogeneous polynomial. The following are equivalent:
    \begin{enumerate}[(i)]
        \item $P$ is stable.
        \item $P$ is hyperbolic with respect to every point in $\R_{>0}^E$.
        \item $P$ is hyperbolic with respect to every point in $\R_{>0}^E$ and every $e\in\R_{\geq0}^E$ with $P(e)\neq0$.
    \end{enumerate}
\end{prop}

\begin{lem}
    Let $P\in\R[x_i\mid i\in E]$ be a homogeneous polynomial and $i\in E$. The following are equivalent:
    \begin{enumerate}[(i)]
        \item $P$ is hyperbolic with respect to $\delta_i$.
        \item $P|_{x_i=1}\in\R[x_i\mid i\in E\smallsetminus\{i\}]$ is a real zero polynomial.
    \end{enumerate}
\end{lem}

\begin{proof}
    We define $Q=P|_{x_i=1}$ and let $d=\deg(P)$.
    Assume that $P$ is hyperbolic with respect to $\delta_i$ and let $v\in\R^{E\smallsetminus\{i\}}$, $\lambda\in\C$ such that $Q(\lambda\cdot v)=0$. We have to show that $\lambda\in\R$. Without loss of generality, we can assume that $\lambda\neq0$. We have
    \begin{equation*}
        0=Q(\lambda\cdot v)=P(\delta_i+\sum_{j\in E\smallsetminus\{i\}}\lambda\cdot v_j\cdot \delta_j)=\lambda^d\cdot P(\lambda^{-1}\cdot\delta_i+\sum_{j\in E\smallsetminus\{i\}}v_j\cdot \delta_j).
    \end{equation*}
    Because $P$ is hyperbolic with respect to $\delta_i$, this shows that $\lambda^{-1}$ and thus $\lambda$ is real.

    Now assume that $Q$ is a real zero polynomial. Let $v\in\R^E$, $\lambda\in\C$ such that $P(\lambda\cdot \delta_i+v)=0$. We have to show that $\lambda\in\R$. We have
    \begin{equation*}
        0=P(\lambda\cdot \delta_i+v)=P((\lambda+v_i)\cdot\delta_i+\sum_{j\in E\smallsetminus\{i\}} v_j\cdot \delta_j).
    \end{equation*}
    Because $v_i\in\R$, it suffices to show that $\mu=\lambda+v_i$ is real. As above, we can assume that $\mu\neq0$. We have
    \begin{equation*}
        0=P(\mu\cdot\delta_i+\sum_{j\in E\smallsetminus\{i\}} v_j\cdot \delta_j)=\mu^{d}\cdot P(\delta_i+\sum_{j\in E\smallsetminus\{i\}} \mu^{-1}\cdot v_j\cdot \delta_j)=Q(\mu^{-1}\cdot\sum_{j\in E\smallsetminus\{i\}} v_j\cdot \delta_j).
    \end{equation*}
    Because $Q$ is a real zero polynomial, this shows that $\mu^{-1}$ and thus $\lambda$ is real.
\end{proof}

\begin{Def}
    Let $P\in\R[x_i\mid i\in E]$ be hyperbolic with respect to $e\in\R^E$. The \emph{hyperbolicity cone} of $P$ at $e$ is defined as
    \begin{equation*}
        \Lambda(P,e)=\{v\in\R^E\mid P(t\cdot e+v)\in\R[t]\textrm{ has only nonnegative zeros}\}.
    \end{equation*}
\end{Def}

Hyperbolicity cones are convex cones \cite{gar}. \Cref{prop:stablehyp} can be rephrased in terms of hypberbolicity cones.

\begin{prop}
   Let $P\in\R[x_i\mid i\in E]$ be a homogeneous polynomial and $e\in\R_{>0}^E$. The following are equivalent:
   \begin{enumerate}[(i)]
        \item $P$ is stable.
       \item $P$ is hyperbolic with respect to $e$ and $\delta_i\in\Lambda(P,e)$ for all $i\in E$.
   \end{enumerate}
\end{prop}

The connection of stable polynomials to M-convex sets is given by the following.

\begin{thm}[{\cite[Theorem~3.2]{BrandenHPP}}]
    Let $P\in\R[x_i\mid i\in E]$ be a homogeneous stable polynomial. Then $\supp(P)$ is M-convex.
\end{thm}

\begin{Def}
    Let $J\subseteq\N_0^E$ be an M-convex set.
    \begin{enumerate}[(a)]
        \item The \emph{generating polynomial} of $J$ is
        \begin{equation*}
            h_J=\sum_{\alpha\in J}\frac{x^\alpha}{\alpha!}\in \R[x_i\mid i\in E]
        \end{equation*}
        where $\alpha!=\prod_{i\in E}\alpha_i!$.
        \item We say that $J$ has the \emph{half-plane property} if its generating polynomial is stable.
        \item We say that $J$ has the \emph{weak half-plane property} if there exists a homogeneous stable polynomial $P\in\R[x_i\mid i\in E]$ such that $J=\supp(P)$.
        \item We say that a polymatroid $r$ on $E$ has the \emph{(weak) half-plane property} if the associated M-convex set $J_r$ has the (weak) half-plane property.
    \end{enumerate}
\end{Def}

In the remaining part, we observe that in certain nice situations, restricting the polymatroid of a homogeneous stable polynomial corresponds to plugging in zeros for some of the variables.

\begin{Def}
    Let $J\subseteq\N_0^E$ be M-convex and let $T\subseteq E$. We say that $J$ is \emph{nondegenerate with respect to $T$} if there exists $\alpha\in J$ such that for all $i\in E\smallsetminus T$ we have~$\alpha_i=0$.
\end{Def}

\begin{lem}\label{lem:mconvrest}
    Let $J\subseteq\N_0^E$ be M-convex and nondegenerate with respect to $T\subseteq E$. Then we have $r_J|_T=r_{J'}$ where
    \begin{equation*}
        J'=\{(\alpha_i)_{i\in T}\mid \alpha\in J\textrm{ such that }\forall i\in E\smallsetminus T: \alpha_i=0\}\subseteq\N_0^T.
    \end{equation*}
\end{lem}

\begin{proof}
    Let $S\subseteq T$. Then
    \begin{equation*}
        r_J|_T(S)=r_J(S)=\max\{\sum_{k\in S}\alpha_k\mid \alpha\in J\}.
    \end{equation*}
    Let $\alpha\in J$ be a point where this maximum is attained such that $\sum_{k\in E\smallsetminus T}\alpha_k$ is minimal. If $\sum_{k\in E\smallsetminus T}\alpha_k=0$, then $\alpha\in J'$ and we have $r_J|_T(S)=r_{J'}(S)$. Thus assume for the sake of a contradiction that $\alpha_i>0$ for some $i\in E\smallsetminus T$. Since $J$ is nondegenerate with respect to $T\subseteq E$, there exists $\beta\in J$ such that for all $k\in E\smallsetminus T$ we have $\beta_k=0$. Because $J$ is M-convex, there exists $j\in T$ such that $\alpha_j<\beta_j$ and $\gamma=\alpha-\delta_i+\delta_j\in J$. Then $\sum_{k\in S}\gamma_k\geq\sum_{k\in S}\alpha_k$ because $i\not\in S$ and $\sum_{k\in E\smallsetminus T}\gamma_k<\sum_{k\in E\smallsetminus T}\alpha_k$ because $i\in E\smallsetminus T$ and $j\in T$. This contradicts our minimality assumption.
\end{proof}

\begin{cor}\label{cor:polyamal}
    Let $E=E_1\cup E_2$ and let $P\in\R[x_i\mid i\in E]$ be a polynomial such that $J=\supp(P)$ is M-convex.
    Finally, for $k=1,2$, we let
    \begin{equation*}
        P_k=P|_{x_i=0\textrm{ for }i\not\in E_k}\in \R[x_i\mid i\in E_k]
    \end{equation*}
    and $J_k=\supp(P_k)$. If $P_1,P_2$ are both not identically zero, then $J_1$ and $J_2$ are M-convex and $r_J$ is an amalgam of $r_{J_1}$ and $r_{J_2}$.
\end{cor}

\begin{proof}
    The condition that $P_k$ is not identically zero is equivalent to $J$ being nondegenerate with respect to $E_k$. Now the claim follows from \Cref{lem:mconvrest}.
\end{proof}

\section{Real zero amalgamation}\label{sec:amal}
In this section we always let $E=E_1\cup E_2$ be a finite set and $E_0=E_1\cap E_2$. We assume that $0\in E_0$. For $k=0,1,2$ we let $0\neq P_k\in\R[x_i\mid i\in E_k]$ be homogeneous and stable such that
\begin{equation*}
    P_0=P_k|_{x_i=0\textrm{ for }i\not\in E_0}
\end{equation*}
for $k=1,2$. Note that this implies that $P_0,P_1$ and $P_2$ all have the same degree $d$. We further denote $J_k=\supp(P_k)$ and $r_k=r_{J_k}$ for $k=0,1,2$. For $k=0,1,2$ we also consider the polynomial $H_k$ obtained from $P_k$ by substituting $x_0+x_i$ for $x_i$ for all~$i\in E_0\smallsetminus\{0\}$.

\begin{lem}\label{lem:hypcone1}
    The polynomial $H_k$ is hyperbolic with respect to $\delta_0$. Its hyperbolicity cone contains $\delta_i$ for all $i\in E_k$ and the point
    \begin{equation*}
        \delta_0-\sum_{0\neq i\in E_0}\delta_i.
    \end{equation*}
\end{lem}

\begin{proof}
    The statement of $H_k$ being hyperbolic with respect to $\delta_0$ is equivalent to $P_k$ being hyperbolic with respect to $\sum_{i\in E_0}\delta_i$. Because of
    \begin{equation*}
        P_k\left(\sum_{i\in E_0}\delta_i\right)=P_0\left(\sum_{i\in E_0}\delta_i\right)\neq0,
    \end{equation*}
    by stability of $P_0$, this follows because $P_k$ is stable. The hyperbolicity cone of $H_k$ containing $\delta_i$, $i\in E_k$, follows from the same statement for $P_k$. The last statement is equivalent to $\delta_0$ being in the hyperbolicity cone of $P_k$.
\end{proof}

\begin{lem}\label{lem:hypcone2}
    Let $H\in\R[x_i\mid i\in E]$ be hyperbolic with respect to $\delta_0$ such that the hyperbolicity cone of
    \begin{equation*}
     H|_{x_i=0\textrm{ for }i\not\in E_k}
    \end{equation*}
    agrees with the hyperbolicity cone of $H_k$ for $k=1,2$. Then the hyperbolicity cone of~$H$ contains $\delta_i$ for all $i\in E$ and the point
    \begin{equation*}
        \delta_0-\sum_{0\neq i\in E_0}\delta_i.
    \end{equation*}
\end{lem}

\begin{proof}
    This follows immediately from \Cref{lem:hypcone1}.
\end{proof}

Now let $E'=E\smallsetminus\{0\}$ and likewise $E_k'=E_k\smallsetminus\{0\}$ for $k=0,1,2$. Then we consider the polynomials
\begin{equation*}
    Q_k=H_k|_{x_0=1}\in\R[x_i\mid i\in E_k']
\end{equation*}
for $k=0,1,2$. These are real zero polynomials with the property that
\begin{equation*}
    Q_0=Q_k|_{x_i=0\textrm{ for }i\not\in E'_0}
\end{equation*}
or $k=1,2$. We have the following result.

\begin{thm}\label{thm:amalmain}
    Assume that there is a real zero polynomial $Q\in\R[x_i\mid i\in E']$ with 
    \begin{equation*}
        Q_k=Q|_{x_i=0\textrm{ for }i\not\in E'_k}
    \end{equation*}
    for $k=1,2$. Then $r_1$ and $r_2$ have an amalgam.
\end{thm}

\begin{proof}
    Let $d'=\deg(Q)$. We define
    \begin{equation*}
        H=x_0^{d'}\cdot Q(\frac{x_i}{x_0}\mid i\in E')\in\R[x_i\mid i\in E].
    \end{equation*}
    Then we have
    \begin{equation*}
        H|_{x_i=0\textrm{ for }i\not\in E'_k}=x_0^{d'-d}\cdot H_k
    \end{equation*}
    for $k=1,2$. Thus by \Cref{lem:hypcone2} the hyperbolicity cone of $H$ contains $\delta_i$, $i\in E$, and the point
    \begin{equation*}
        \delta_0-\sum_{0\neq i\in E_0}\delta_i.
    \end{equation*}
    This implies that the polynomial $P$ obtained from $H$ by substituting $x_i-x_0$ for $x_i$ for all $i\in E_0\smallsetminus\{0\}$ is stable. We have
    \begin{equation*}
        P|_{x_i=0\textrm{ for }i\not\in E_k}=x_0^{d'-d}\cdot P_k
    \end{equation*}
    for $k=1,2$. Finally, we let $P'$ be the polynomial obtained from $P$ by dropping all monomials that are not divisible by $x_0^{d'-d}$ and dividing the result by $x_0^{d'-d}$. Then
    \begin{equation*}
        P'|_{x_i=0\textrm{ for }i\not\in E_k}=P_k
    \end{equation*}
    for $k=1,2$. The support of $P'$ is M-convex because it agrees with the support of the stable polynomial $\frac{\partial^{d'-d}}{\partial x_0}P$. Now the claim follows from \Cref{cor:polyamal}.
\end{proof}

Now we are ready to disprove the real zero amalgamation conjecture from \cite{rzamalgamation}.

\begin{con}[{\cite[Conjecture~7.6]{rzamalgamation}}]\label{con:weakrzamal}
    Let $E'=E_1'\cup E_2'$ be a finite set such that $E_0'=E_1'\cap E_2'$ has two elements. For $k=1,2$ let $Q_k\in\R[x_i\mid i\in E_k']$ be a real zero polynomial.
    If
    \begin{equation*}
        Q_1|_{x_i=0\textrm{ for }i\not\in E_0'}=Q_2|_{x_i=0\textrm{ for }i\not\in E_0'},
    \end{equation*}
    then there is a real zero polynomial $Q\in\R[x_i\mid i\in E']$ such that 
    \begin{equation*}
     Q_k=Q|_{x_i=0\textrm{ for }i\not\in E_k'}   
    \end{equation*}
     for $k=1,2$.
\end{con}

\Cref{thm:amalmain} shows that counterexamples to \Cref{con:weakrzamal} can be constructed from polymatroids with the half-plane property that do not have an amalgam. In the next subsection we construct such a counterexample with $|E_0|=3$.

\begin{figure}
        \begin{tikzpicture}[line cap=round,line join=round,>=triangle 45,x=.75cm,y=0.75cm]
\clip(-4,-4) rectangle (4,4);
\draw [line width=2pt] (0,3)-- (-2.6,-1.5);
\draw [line width=2pt] (0,3)-- (2.6,-1.5);
\draw [line width=2pt] (0,3)-- (0,-1.5);
\fill [color=black] (0,3) circle (4.5pt);
\draw[color=black] (0,3.5) node {\textbf{3}};
\fill [color=black] (0,0) circle (4.5pt);
\draw[color=black] (0.45,0.25) node {\textbf{1}};
\fill [color=black] (-2.6,-1.5) circle (4.5pt);
\draw[color=black] (-3,-1.75) node {\textbf{0$'$}};
\fill [color=black] (2.6,-1.5) circle (4.5pt);
\draw[color=black] (3,-1.75) node {\textbf{2$'$}};
\fill [color=black] (-1.3,0.75) circle (4.5pt);
\draw[color=black] (-1.75,1) node {\textbf{0}};
\fill [color=black] (1.3,0.75) circle (4.5pt);
\draw[color=black] (1.75,1) node {\textbf{2}};
\fill [color=black] (0,-1.5) circle (4.5pt);
\draw[color=black] (0,-2) node {\textbf{1$'$}};
\end{tikzpicture}
\begin{tikzpicture}[line cap=round,line join=round,>=triangle 45,x=0.75cm,y=0.75cm]
\clip(-4,-4) rectangle (4,4);
\draw [line width=2pt] (0,3)-- (-2.6,-1.5);
\draw [line width=2pt] (0,3)-- (2.6,-1.5);
\fill [color=black] (0,3) circle (4.5pt);
\draw[color=black] (0,3.5) node {\textbf{4}};
\fill [color=black] (0,0) circle (4.5pt);
\draw[color=black] (0.45,0.25) node {\textbf{1}};
\fill [color=black] (-2.6,-1.5) circle (4.5pt);
\draw[color=black] (-3,-1.75) node {\textbf{0$'$}};
\fill [color=black] (2.6,-1.5) circle (4.5pt);
\draw[color=black] (3,-1.75) node {\textbf{2$'$}};
\fill [color=black] (-1.3,0.75) circle (4.5pt);
\draw[color=black] (-1.75,1) node {\textbf{0}};
\fill [color=black] (1.3,0.75) circle (4.5pt);
\draw[color=black] (1.75,1) node {\textbf{2}};
\fill [color=black] (0,-1.5) circle (4.5pt);
\draw[color=black] (0,-2) node {\textbf{1$'$}};
\end{tikzpicture}
        \caption{The matroids $F_7^{-4}$ (left) and $F_7^{-5}$ (right).}
        \label{fig:f7relax}
 \end{figure}

\subsection{An example with \texorpdfstring{$|E_0|=3$}{E0=3}}
We consider the matroids $F_7^{-4}$ and $F_7^{-5}$ defined as in \Cref{fig:f7relax}. Here, for $F_7^{-4}$, every 3-element subset of $\{0,0',1,1',2,2',3\}$ forms a basis, except for the subsets $\{0',0,3\}$, $\{1',1,3\}$ and $\{2',2,3\}$. For $F_7^{-5}$, similarly, all 3-element subsets of $\{0,0',1,1',2,2',4\}$ are bases except for $\{0',0,4\}$ and $\{2',2,4\}$. These matroids both have the half-plane property, see \cite[\S A.2.2]{HPPlong} for $F_7^{-5}$ and \cite{wagnerwei} for~$F_7^{-4}$. Denote by $e_3$ the third elementary symmetric polynomials in 7 variables. This means that their generating polynomials
\begin{equation*}
    h_{F_7^{-4}}=e_3(x_0,x_{0'},x_1,x_{1'},x_2,x_{2'},x_3)-x_{0'}x_0x_3-x_{1'}x_1x_3-x_{2'}x_2x_3
\end{equation*}
and
\begin{equation*}
    h_{F_7^{-5}}=e_3(x_0,x_{0'},x_1,x_{1'},x_2,x_{2'},x_4)-x_{0'}x_0x_4-x_{2'}x_2x_4
\end{equation*}
are stable. We let $P_1$ and $P_2$ be the polynomials obtained from $h_{F_7^{-4}}$ and $h_{F_7^{-5}}$, respectively, by setting $x_{0'}=x_0$, $x_{1'}=x_1$ and $x_{2'}=x_2$. We have
\begin{align*}
    P_1&=4x_1x_2x_3+2x_1x_2^2+2x_1^2x_2+4x_0x_2x_3+2x_0x_2^2+4x_0x_1x_3+8x_0x_1x_2\\
    &+2x_0x_1^2+2x_0^2x_2+2x_0^2x_1
\end{align*}
and
\begin{align*}
    P_2&=4x_1x_2x_4+2x_1x_2^2+x_1^2x_4+2x_1^2x_2+4x_0x_2x_4+2x_0x_2^2+4x_0x_1x_4+8x_0x_1x_2\\
    &+2x_0x_1^2+2x_0^2x_2+2x_0^2x_1.
\end{align*}
By construction we have
\begin{equation*}
    P_1|_{x_3=0}=P_2|_{x_4=0}.
\end{equation*}
Hence $E_1=\{0,1,2,3\}$ and $E_2=\{0,1,2,4\}$. Note that the support of $P_0\coloneqq P_1|_{x_3=0}$ is the M-convex set from \Cref{ex:uniform} for $n=3$. We will show that the polymatroids $r_1$ and $r_2$ corresponding to $P_1$ and $P_2$ do not have an amalgam so that the real zero polynomials $Q_1$ and $Q_2$ obtained from $P_1$ and $P_2$ constitute a counterexample to \Cref{con:weakrzamal}. In fact, we will even prove that for every $m\in\N$ the polymatroids $m\cdot r_1$ and $m\cdot r_2$ do not have an amalgam which shows that $Q_1^m$ and $Q_2^m$ also do not satisfy the conclusion of \Cref{con:weakrzamal} for any $m\in\N$.

\begin{thm}
    The polymatroids $m\cdot r_1$ and $m\cdot r_2$ do not have an amalgam.
\end{thm}

\begin{proof}
    We proceed as in the proof of \cite[Theorem~2]{sticky}.
    Assume for the sake of a contradiction that the polymatroid $r$ on $E=\{0,1,2,3,4\}$ is an amalgam of $m\cdot r_1$ and $m\cdot r_2$. We have
    \begin{equation*}
        r(\{0\})+r(\{0,3\})\leq r(\{0\})+r(\{0,3,4\})\leq r(\{0,3\})+r(\{0,4\}).
    \end{equation*}
    By definition of $r_1$ and $r_2$ this implies that
    \begin{equation*}
        2m+2m\leq 2m+r(\{0,3,4\})\leq 2m+2m.
    \end{equation*}
    Thus we have $r(\{0,3,4\})=2m$. Likewise one shows that $r(\{2,3,4\})=2m$. Furthermore, we have that
    \begin{equation*}
      4m\leq r(\{3,4\})+3m \leq r(\{3,4\})+r(\{0,2,3,4\})\leq r(\{0,3,4\})+r(\{2,3,4\})=4m.
    \end{equation*}
    This shows $r(\{3,4\})=m$. Finally, we have
    \begin{equation*}
      m+2m\leq  r(\{3\})+r(\{1,3,4\}) \leq r(\{1,3\})+r(\{3,4\}) =2m+m
    \end{equation*}
    which implies $r(\{1,3,4\})=2m$ contradicting
    \begin{equation*}
        3m=r(\{1,4\})\leq r(\{1,3,4\}).
     \end{equation*}
     Hence the polymatroids $m\cdot r_1$ and $m\cdot r_2$ do not have an amalagam.
\end{proof}

\section{Weak half-plane property}\label{sec:whpp}
Let $E$ always denote a finite set and $M$ a matroid on $E$ with set of bases $\cB$.

\begin{Def}[\cite{BrandenDLeon}]
    We say that four bases $B_1,B_2,B_3,B_4\in\cB$ form a \emph{degenerate quadrangle} of $M$ if there exists $S\subseteq E$ and pairwise different $i,j,k,l\notin S$ such that 
    \begin{equation*}
      (B_1,B_2,B_3,B_4)=(S\cup\{i,k\},S\cup\{i,l\},S\cup\{j,l\},S\cup\{j,k\})  
    \end{equation*}
    and if at most one of $S\cup\{i,j\}$ and $S\cup\{k,l\}$ is a bases of $M$. 
\end{Def}

The following theorem was used in \cite{BrandenDLeon} to reduce the number of possible parameters when searching for a stable polynomial with support $M$.

\begin{thm}[\cite{BrandenHPP}]\label{thm:degQuadCoeffs}
    For every basis $B\in\cB$ let $0\neq a_B\in\R$ be such that the multiaffine and homogeneous polynomial
    \begin{equation*}
        P=\sum_{B\in\cB}a_{B}\cdot\prod_{i\in B}x_i\in\R[x_i\mid i\in E]
    \end{equation*}
    is stable. If $B_1,B_2,B_3,B_4$ form a degenerate quadrangle of $M$, then
    \begin{equation}\label{eq:deg4}
      a_{B_1}a_{B_3}=a_{B_2}a_{B_4}.  
    \end{equation}
\end{thm}
 Letting $b_B\coloneqq \log(|a_B|)$ for all $B\in\cB$ we obtain from \Cref{eq:deg4} linear equations
 \begin{equation*}
  b_{B_1}+b_{B_3}-b_{B_2}-b_{B_4}=0   
 \end{equation*}
for all degenerate quadrangles $B_1,B_2,B_3,B_4$ of $M$. We denote by $V_M\subseteq\R^{\cB}$ the linear space cut out by all such equations. By \cite[Lemma~2.6]{BrandenDLeon} the vector space
\begin{equation*}
  W_M\coloneqq \{(\sum_{i\in B}v_i)_{B\in\cB}\mid v\in\R^E\}  
\end{equation*}
is an $(|E|-z+1)$-dimensional linear subspace of $V_M$ where $z$ is the number of connected components of $M$.

\begin{lem}\label{lem:rescale}
    Let $U_M$ be a linear complement of $W_M$ in $V_M$. If $M$ has the weak half-plane property, then there exists a vector $b\in U_M$ such that
    \begin{equation*}
        \sum_{B\in\cB}\exp(b_{B})\cdot\prod_{i\in B}x_i
    \end{equation*}
    is stable.
\end{lem}

\begin{proof}
    This follows in the same way as \cite[Theorem~2.3]{BrandenDLeon}: Scaling the variables by $x_i\mapsto\exp(v_i)x_i$ corresponds, after taking logarithms of the coefficients, to shifting by the corresponding vector from $W_M$.
\end{proof}

For choosing a linear complement of $W_M$ in $V_M$ in a nice way, the following lemma might be useful.

\begin{lem}\label{lem:notreg}
    Let $M$ be represented by a matrix $A\in\R^{d\times |E|}$ of rank $d$. For $B\in\cB$ we denote by $A[B]$ the corresponding $d\times d$ submatrix. Then 
    \begin{equation*}
        u(A)\coloneqq (\log|\det(A[B])|)_{B\in\cB}\in V_M.
    \end{equation*}
    Furthermore, if $M$ is not regular, then $u(A)\not\in W_M$.
\end{lem}

\begin{proof}
    By \cite[Theorem~8.1]{HPPlong} the polynomial
    \begin{equation*}
        \sum_{B\in\cB}\det(A[B])^2\cdot\prod_{i\in B}x_i
    \end{equation*}
    is stable. This, together with \Cref{thm:degQuadCoeffs}, proves the first claim. Now assume that there exists $v\in\R^E$ such that $u(A)=(\sum_{i\in B}v_i)_{B\in\cB}$. Scaling the $i$th column of $A$ by $\exp(-v_i)$ for all $i\in E$, we obtain a matrix $A'$ representing $M$ all of whose maximal minors are in $\{-1,0,1\}$. After multiplication of $A'$ from the left by a suitable invertible matrix we can additionally assume that $A'[B_0]$ is the identity matrix for some $B_0\in\cB$. Then $A'$ is a totally unimodular matrix representing $M$ which shows that $M$ is regular.
\end{proof}

\begin{ex}\label{ex:p8}
    In this example, we consider the rank 4 matroid $M=P_8$ on 8 elements which is represented by the real matrix
    \[A\coloneqq \begin{pmatrix}
        1&0&0&0&0&1&1&2\\
        0&1&0&0&1&0&1&1\\
        0&0&1&0&1&1&0&1\\
        0&0&0&1&2&1&1&0
    \end{pmatrix}\]
    whose columns we label by $0,\ldots,7$.
    This matroid is not regular \cite[\S A4]{HPPlong} and therefore $u(A)$ is in $V_M$ but not in $W_M$ by \Cref{lem:notreg}. Using the {\texttt{Macaulay2}} \cite{M2} package ``Matroids'' \cite{MatroidsArticle} we compute that $\dim(V_M)=9$. Because $M$ is connected, this implies that the span of $u(A)$ is a linear complement of $W_M$ in $V_M$. Note that $P_8$ has the weak half-plane property because it is representable over $\R$. 
\end{ex}

Recall that if $X$ is a circuit-hyperplane of $M$, then $\cB\cup\{X\}$ is the set of bases of a matroid $M'$ \cite[Theorem~1.5.14]{oxley}. Then $M'$ is called a \emph{relaxation} of $M$.

\begin{lem}\label{lem:relax}
    Let $M'$ be a relaxation of $M$. Then every degenerate quadrangle of $M'$ is a degenerate quadrangle of $M$.
\end{lem}

\begin{proof}
    We proceed as in the proof of \cite[Lemma~3.43]{masterleon}.
    Denote by $X$ the circuit-hyperplane of $M$ such that $M'$ is the relaxation of $M$ by $X$.
    We show that for all $x\in X$ and $y\in E\smallsetminus X$ the set
    \[(X\smallsetminus\{x\})\cup\{y\}\]
    forms a basis of $M$ (and therefore of $M'$). Then it immediately follows that $X$ cannot be contained in some degenerate quadrangle of $M'$.
    Since $X$ is a circuit, $X\smallsetminus\{x\}$ still has rank $\rank(M)-1=\rank(M')-1$. Then $(X\smallsetminus\{x\})\cup\{y\}$ has rank $\rank(M')$ because $X$ is closed and $y\notin X$. Thus $(X\smallsetminus\{x\})\cup\{y\}$ is a basis of $M'$.
\end{proof}

\begin{rem}\label{rem:embed}
    Assume that $\rank(M)\geq2$.
    Let $M'$ be the relaxation of $M$ by the circuit-hyperplane $X$ and denote $\cB'=\cB\cup\{X\}$. \Cref{lem:relax} implies that the map $\R^{\cB}\to\R^{\cB'}$, that sends $v\in\R^{\cB}$ to the vector $w$ with $w_{X}=0$ and $w_B=v_B$ for all $B\in\cB$, maps $V_M$ to $V_{M'}$. We thus obtain a natural embedding
    \begin{equation*}
        \iota\colon V_M\hookrightarrow V_{M'}.
    \end{equation*}
    Furthermore, we have $\delta_X\in V_{M'}$. If $U_M$ is a linear complement of $W_M$ in $V_M$, then
    \begin{equation*}
        \left(\iota(U_M)\oplus\R\cdot\delta_X\right)\cap W_{M'}=\{0\}.
    \end{equation*}
    However, in general $\iota(U_M)$ and $\delta_X$ do not span a linear complement of $W_{M'}$, see for example \cite[Table~1]{BrandenDLeon}.
\end{rem}

\begin{ex}\label{ex:f7}
    As an illustration, we recall \cite[Example~4.1]{BrandenDLeon}. The non-Fano matroid $M'=F_7^{-}$ is represented by the real matrix
    \[A=\begin{pmatrix}
        1&1&0&0&0&1&1\\
        0&1&1&1&0&0&1\\
        0&0&0&1&1&1&1
    \end{pmatrix}\]
    whose columns we label by $1,\ldots,7$. This matroid is not regular \cite[\S A.2.2]{HPPlong} and therefore $u(A)$ is in $V_{M'}$ but not in $W_{M'}$ by \Cref{lem:notreg}. On the other hand, the non-Fano matroid is a relaxation of the Fano matroid $F_7$ by the circuit hyperplane $\{2,4,6\}$. Because $\dim(V_{M'}/W_{M'})=1$ by \cite{BrandenDLeon}, it follows from \Cref{rem:embed} that $u(A)+W_{M'}$ must contain a scalar multiple of $\delta_{\{2,4,6\}}$. Indeed, one has that $u(A)=\log(2)\cdot\delta_{\{2,4,6\}}$. It was further noted in \cite[Example~4.1]{BrandenDLeon} that by \cite[Example~11.5]{HPPlong} the polynomial
    \begin{equation*}
        h_{F_7}+\mu x_2x_4x_6
    \end{equation*}
    is stable only for $\mu=4$. As a side note we would like to mention that this implies in particular that
    \begin{equation*}
        \sum_{S\in\binom{[7]}{3}}|\det(A[S])|\cdot\prod_{i\in S}x_i
    \end{equation*}
    is not stable, giving a negative answer to the question raised in \cite[Remark~4.2]{purbhoo}.
\end{ex}

Now we are ready to disprove the following conjecture by Br{\"a}nd{\'e}n--{D}'Le{\'o}n.

\begin{con}[{\cite[Conjecture~4.2]{BrandenDLeon}}]\label{con:relax}
    Suppose that $M$ has the weak half-plane property. Then so does any relaxation of $M$.
\end{con}

Our counterexample is a suitable relaxation of $P_8$.

\begin{ex}\label{ex:counterrelax}
    Consider again the matroid $P_8$ from \Cref{ex:p8}. The set $X\coloneqq \{3,5,6,7\}$ is a circuit-hyperplane of $P_8$. Following \cite{ninematroids} we denote by $P_1$ the relaxation of $P_8$ by $X$. We will prove that $P_1$ does not have the weak half-plane property, although it is a relaxation of $P_8$.    
    Using the {\texttt{Macaulay2}} \cite{M2} package ``Matroids'' \cite{MatroidsArticle} we compute that $\dim(V_{P_1})=10$. Thus by \Cref{ex:p8} and \Cref{rem:embed} the vectors $\delta_X$ and $\iota(u(A))$ span a linear complement of $W_{P_1}$ in $V_{P_1}$. We define $v=\frac{1}{\log(2)}\iota(u(A))$ --- this is just for convenience to get a vector with entries in $\{0,1,2\}$. If $P_1$ has the weak half-plane property, then by \Cref{lem:rescale} there are $a,b>0$ such that
    \begin{equation*}
        F_{a,b}=\sum_{B\text{ basis of }P_8}b^{v_B}x^B+ax_3x_5x_6x_7
    \end{equation*}
    is stable. Now consider the $(0,1)$-th Rayleigh difference
    \begin{equation*}
      \Delta_{0,1}F_{a,b}\coloneqq \partial_{x_0}F_{a,b}\partial_{x_1}F_{a,b}-F_{a,b}\partial_{x_0}\partial_{x_1}F_{a,b}.  
    \end{equation*}
    By \cite[Theorem~5.10]{BrandenHPP} the Rayleigh difference $\Delta_{0,1}F_{a,b}$ is globally nonnegative if $F_{a,b}$ is stable. Plugging in $(1,1,t,-1,-1,t)$ for $(x_2,\ldots,x_8)$ yields
    \begin{equation*}
      (\Delta_{0,1}F_{a,b})(1,1,t,-1,-1,t)=-abt^3+(-ab-4b^2+2a+12b+16)t^2+at.  
    \end{equation*}
     Because $a,b>0$, this takes negative values for large enough $t$. Therefore, for no choice of $a,b>0$ the polynomial $F_{a,b}$ is stable.
    
\end{ex}

\begin{rem}
    According to \cite[Proposition~4]{ninematroids} the matroid $P_1$ is not representable over any field. It does not have the half-plane property because it has the matroid $F_7^{-3}$ as a minor which does not have the half-plane property by \cite[Example~11.7]{HPPlong}. Moreover, it does not seem to be representable even in the more general context studied in \cite{skewpartial}. This made $P_1$ a good candidate for being a counterexample to \Cref{con:relax} as representations over certain algebras (that are not necessarily fields) sometimes can still be used to prove the weak half-plane property. See for instance \cite{NonPappus} where the weak half-plane property was proved for the non-Pappus and  the non-Desargues matroid.
\end{rem}

\begin{rem}\label{rem:subd}
    We use the notation as in \Cref{ex:counterrelax}.
    The polynomial $F_{0,1}\coloneqq \lim_{a\to0}F_{a,1}$ is the basis generating polynomial of $P_8$. The polynomial $F_{0,0}\coloneqq \lim_{a,b\to0}F_{a,b}$ is the basis generating polynomial of the graphical matroid $M(G_1)$ where $G_1$ is depicted in \Cref{fig:graphs} on the left. In total, the regular subdivision of the matroid polytope of $P_8$ defined by $v_B$ has six maximal cells, all of which are matroid polytopes themselves. Four of the corresponding matroids are isomorphic to $M(G_1)$, the two remaining ones are isomorphic to $M(G_2)$ where $G_2$ is the graph on the right of \Cref{fig:graphs}.
    Unfortunately, this knowledge did not help us representing $F_{a,b}$ in a simple way.
\end{rem}

\begin{figure}
\begin{tikzpicture}[line cap=round,line join=round,>=triangle 45,x=0.75cm,y=0.75cm]
\clip(-4,-4) rectangle (4,4);
\draw [line width=2pt] (0,0)-- (-2.6,-1.5);
\draw [line width=2pt] (0,0)-- (2.6,-1.5);
\draw [line width=2pt] (0,3)-- (-2.6,-1.5);
\draw [line width=2pt] (0,3)-- (2.6,-1.5);
\draw [shift={(-4.49,1.5)},line width=2pt]  plot[domain=-0.32:0.32,variable=\t]({1*4.73*cos(\t r)+0*4.73*sin(\t r)},{0*4.73*cos(\t r)+1*4.73*sin(\t r)});
\draw [shift={(4.5,1.5)},line width=2pt]  plot[domain=2.82:3.46,variable=\t]({1*4.74*cos(\t r)+0*4.74*sin(\t r)},{0*4.74*cos(\t r)+1*4.74*sin(\t r)});
\draw [line width=2pt] (-2.6,-1.5)-- (0,-1.5);
\draw [line width=2pt] (0,-1.5)-- (2.6,-1.5);
\fill [color=black] (0,0) circle (4.5pt);
\fill [color=black] (0,3) circle (4.5pt);
\fill [color=black] (-2.6,-1.5) circle (4.5pt);
\fill [color=black] (2.6,-1.5) circle (4.5pt);
\draw[color=black] (-1.3,-0.5) node {\textbf{6}};
\draw[color=black] (1.3,-0.5) node {\textbf{1}};
\draw[color=black] (-1.5,0.9) node {\textbf{5}};
\draw[color=black] (1.5,0.9) node {\textbf{2}};
\fill [color=black] (0,-1.5) circle (4.5pt);
\draw[color=black] (0.5,1.4) node {\textbf{4}};
\draw[color=black] (-0.5,1.4) node {\textbf{7}};
\draw[color=black] (-1.4,-1.7) node {\textbf{0}};
\draw[color=black] (1.4,-1.7) node {\textbf{3}};
\end{tikzpicture}
\begin{tikzpicture}[line cap=round,line join=round,>=triangle 45,x=0.75cm,y=0.75cm]
\clip(-4,-4) rectangle (4,4);
\draw [line width=2pt] (0,3)-- (-3,0);
\draw [line width=2pt] (-3,0)-- (0,-3);
\draw [line width=2pt] (0,-3)-- (3,0);
\draw [line width=2pt] (3,0)-- (0,3);
\draw [shift={(-1.5,4.5)},line width=2pt]  plot[domain=4.39:5.03,variable=\t]({1*4.75*cos(\t r)+0*4.75*sin(\t r)},{0*4.75*cos(\t r)+1*4.75*sin(\t r)});
\draw [shift={(-1.5,-4.5)},line width=2pt]  plot[domain=1.25:1.89,variable=\t]({1*4.74*cos(\t r)+0*4.74*sin(\t r)},{0*4.74*cos(\t r)+1*4.74*sin(\t r)});
\draw [shift={(1.5,-4.5)},line width=2pt]  plot[domain=1.25:1.89,variable=\t]({1*4.74*cos(\t r)+0*4.74*sin(\t r)},{0*4.74*cos(\t r)+1*4.74*sin(\t r)});
\draw [shift={(1.5,4.5)},line width=2pt]  plot[domain=4.39:5.03,variable=\t]({1*4.74*cos(\t r)+0*4.74*sin(\t r)},{0*4.74*cos(\t r)+1*4.74*sin(\t r)});
\fill [color=black] (0,3) circle (4.5pt);
\fill [color=black] (-3,0) circle (4.5pt);
\fill [color=black] (0,-3) circle (4.5pt);
\fill [color=black] (3,0) circle (4.5pt);
\fill [color=black] (0,0) circle (4.5pt);
\end{tikzpicture}
        \caption{The graphs $G_1$ (left) and $G_2$ (right) from \Cref{rem:subd}.}
        \label{fig:graphs}
 \end{figure}

\section{Quaternionic unimodular matroids}\label{sec:qu}
We recall the definition of quaternionic unimodular (QU) matroids. To this end, let~$\HH$ denote the skew field of quaternions.

\begin{Def}
    Let $E$ be a finite set and $C\subseteq\HH^E$ a submodule of the free left $\HH$-module $\HH^E$. A nonzero element $x\in C$ is called an \emph{elementary chain} of $C$ if $C$ does not contain a nonzero element whose support is strictly contained in the support of~$x$. The submodule $C$ is called \emph{unimodular} if for every elementary chain $x$ of $C$ all nonzero entries of $x$ have the same norm.
\end{Def}

\begin{thm}[{\cite[Theorem~3.7]{skewpartial}}]
    Let $E$ be a finite set and $C\subseteq\HH^E$ unimodular. The set of supports of elementary chains in $C$ is the set of cocircuits of a matroid $M(C)$ on $E$.
\end{thm}

\begin{Def}\label{def:qu}
    A matroid on a finite set $E$ of the form $M(C)$ for some unimodular $C\subseteq\HH^E$ is called \emph{quaternionic unimodular (QU)}.
\end{Def}

The goal of this section is to prove that every QU matroid has the half-plane property. This has been conjectured in \cite[Conjecture~6.9]{skewpartial}.


\begin{Def}[{\cite[page 219]{skewpartial}}]
    Denote by $\varphi\colon\HH\to\C^{2\times 2}$ the map 
\begin{equation*}
    \varphi(a+bi+cj+dk)=
    \begin{pmatrix}
        a+bi& c+di\\
        -c+di& a-bi
    \end{pmatrix}.
\end{equation*}
    We extend $\varphi$ to matrices by applying $\varphi$ entry-wise. Thus for a matrix $A\in\HH^{n\times n}$ we obtain the $2n\times2n$ complex matrix $\varphi(A)$. We define
    \begin{equation*}
      \delta(A)\coloneqq \sqrt{|\det(\varphi(A))|}.  
    \end{equation*}
\end{Def}

\begin{rem}\label{rem:deltaHom}We collect some basic properties of $\delta$.
\begin{enumerate}
    \item For columns $a_1,\ldots,a_n\in\HH^n$ and $\lambda\in\R$, we have
    \begin{equation*}
      \delta(\lambda a_1\ a_2\ldots a_n)=|\lambda|\cdot\delta(a_1\ldots a_n).  
    \end{equation*}
    This follows directly from the definition.
    \item For matrices $A,B\in\HH^{n\times n}$ we have $\delta(AB)=\delta(A)\delta(B)$ and $\delta(A)=\delta(A^*)$. This is \cite[Lemma~5.2]{skewpartial}.
\end{enumerate}
\end{rem}

For a matrix $A\in\HH^{m\times n}$ with $m\leq n$ and $B\subseteq[n]$ of size $m$ we denote by $A[B]$ the $m\times m$ submatrix of $A$ consisting of the columns indexed by $B$.
The following is a version of the Cauchy--Binet theorem over $\HH$.

\begin{prop}[{\cite[Theorem 5.1]{skewpartial}}]\label{prop:CauchyBinetQuat}
    Let $A\in\HH^{m\times n}$ be a matrix over the quaternions and $m\leq n$. Then
    \begin{equation*}
        \delta(AA^*)=\sum_{\substack{B\subseteq[n]\\|B|=m}}\delta(A[B]A[B]^*).
    \end{equation*}
\end{prop}

\begin{lem}\label{lem:basisIffDelta1}
    Let $E$ be a finite set and $C\subseteq\HH^E$ unimodular. Let $d$ be the rank of the matroid $M=M(C)$.
    There is a $d\times |E|$ matrix $A$ over $\HH$ whose rows form a basis of~$C$ such that for all $B\subseteq E$ of size $d$ we have $\delta(A[B])=1$ if $B$ is a basis of $M$ and~$\delta(A[B])=0$ otherwise.
\end{lem}

\begin{proof}
 Let $A$ be a matrix over $\HH$ whose rows form a basis of $C$. By \cite[Lemma~3.14]{skewpartial} it has $d$ rows, so $A\in\HH^{d\times |E|}$. After multiplying $A$ from the left by an invertible $d\times d$ matrix over $\HH$, we can assume by \cite[Corollary~3.26]{skewpartial} that $A$ is a \emph{strong QU matrix} in the sense of \cite[Definition~3.23]{skewpartial}. Now the statement of the lemma follows from {\cite[Claim 5.4.1]{skewpartial}}.
\end{proof}

\begin{thm}[{\cite[Conjecture~6.9]{skewpartial}}]
    Every QU matroid has the half-plane property.
\end{thm}
\begin{proof}
  Let $E=\{1,\ldots,m\}$ and $C\subseteq\HH^m$ be unimodular. Denote by $M=M(C)$ the associated QU matroid and let $A\in\HH^{d\times m}$ be a matrix as in \Cref{lem:basisIffDelta1}. We have to show that the bases generating polynomial $h_M$ of $M$ is stable. This is equivalent to $h_M^2$ being stable. We prove this by showing that $h_M^2$ agrees with the stable polynomial
  \begin{equation*}
    \det\left(\varphi(A)\begin{pmatrix}x_1&&&&\\&x_1&&&\\&&\ddots&&\\&&&x_m&\\&&&&x_m\end{pmatrix}\varphi(A)^*\right).  
  \end{equation*}
    It suffices to show that these two polynomials agree on the positive orthant. We denote by $a_1,\ldots,a_m\in\HH^d$ the columns of $A$. Let $x\in\R_{>0}^m$ and write $x=(x_1^2,\ldots,x_m^2)$ for $x_1,\ldots,x_m\in\R_{>0}$. Then we have
    \allowdisplaybreaks\begin{align}
        h_M(x)^2&=\left(\sum_{\substack{B\subseteq E\\|B|=d}}\delta(A[B]A[B]^*)x^B\right)^2\\
        &=\left(\sum_{\substack{B\subseteq E\\|B|=d}}\delta((x_1a_1\ldots x_ma_m)[B](x_1a_1\ldots x_ma_m)[B]^*)\right)^2\\
        &=\delta((x_1a_1\ldots x_ma_m)(x_1a_1\ldots x_ma_m)^*)^2\\
        &=|\det(\varphi(x_1a_1\ldots x_ma_m)\varphi(x_1a_1\ldots x_ma_m)^*)|\\
        &=|\det\left(\varphi(A)\begin{pmatrix}x_1^2&&&&\\&x_1^2&&&\\
        &&\ddots&&\\&&&x_m^2&\\&&&&x_m^2\end{pmatrix}\varphi(A)^*\right)|\\
        &=\det\left(\varphi(A)\begin{pmatrix}x_1^2&&&&\\&x_1^2&&&\\
        &&\ddots&&\\&&&x_m^2&\\&&&&x_m^2\end{pmatrix}\varphi(A)^*\right).
    \end{align}
    Here we have equality in (2) by \Cref{rem:deltaHom} and \Cref{lem:basisIffDelta1}. (3) holds  by \Cref{rem:deltaHom}. For (4) we use \Cref{prop:CauchyBinetQuat} and (5) follows from the definition of $\delta$. (6) and (7) are obvious.
\end{proof}

